% Auto-generated by generate_tex_fragments.py -- do not hand-edit.
\begin{longtable}{@{}L{2.0cm}L{2.0cm}L{2.1cm}L{6.2cm}L{1.9cm}@{}}
\caption{P2: A routine private channel makes materially relevant concerns more visible. Record, design, population/setting, finding as charted (verbatim from the library record), and direction relative to the proposition. n=18 linked records (kg/graph.json edges).}\label{tab:p2}\\
\toprule
\textbf{Record} & \textbf{Design} & \textbf{Population\slash{}setting} & \textbf{Finding (as charted)} & \textbf{Direction}\\
\midrule
\endfirsthead
\multicolumn{5}{l}{\textit{P2 continued}}\\
\toprule
\textbf{Record} & \textbf{Design} & \textbf{Population\slash{}setting} & \textbf{Finding (as charted)} & \textbf{Direction}\\
\midrule
\endhead
\bottomrule
\endfoot
Abramsky T 2014 \cite{ref24} & Pair-matched cluster randomized controlled trial (8 \seqsplit{communities}) & Community members aged 18-49, Kampala, Uganda; baseline n=1,583, follow-up n=2,532 & SASA! was associated with \seqsplit{significantly} lower social acceptance of IPV among women (aRR 0.54, 95\% CI 0.38-0.79) and 52\% lower past-year physical IPV among women (0.48, 95\% CI 0.16-1.39, crossing 1), plus higher likelihood of women receiving supportive community responses to disclosed violence. & Supports\\
\\[3pt]
Azalee M 2025 \cite{srp40116416} & \seqsplit{Qualitative} in-depth interviews (n=6) & Health care workers who were IPV survivors, tertiary hospital, Malaysia & Six HCW-survivors used varied coping strategies (seeking help, \seqsplit{spirituality}, avoidance, self-harm, staying\slash{}leaving); survivors relied mainly on coworkers and sought formal help only when situations became critical, hindered by \seqsplit{misconceptions} about IPV, privacy concerns and fear of workplace gossip. & Supports\\
\\[3pt]
Chen 2013 \cite{ref94} & clinical practice review & primary\slash{}family care settings & Describes RADAR (Remember to ask routinely, Ask directly in private, Document, Assess safety, Review options) as a structured clinical tool for \seqsplit{identifying} and responding to IPV, \seqsplit{emphasizing} private \seqsplit{questioning}, safety assessment (weapons, children, \seqsplit{perpetrator} risk), and advocacy-based referral requiring trained links to community resources. & Supports\\
\\[3pt]
Decker 2017 \cite{ref95} & quasi-\seqsplit{experimental} single-group pretest-posttest + \seqsplit{qualitative} interviews (survey n=132 baseline, n=68 3-month retention; interviews n=35) & women aged 18+ at two US family planning clinics & A brief, trauma-informed, universal IPV\slash{}\seqsplit{reproductive}-coercion assessment and education \seqsplit{intervention} was \seqsplit{implemented} with 65\% of women receiving at least one \seqsplit{intervention} element; the mixed-methods evaluation describes uptake and provider\slash{}patient experience of the assessment process. & Supports\\
\\[3pt]
Ford-Gilboe 2016 \cite{ref97} & instrument-\seqsplit{development} study using pooled secondary data from 5 Canadian studies, n=6,278 adult women; \seqsplit{international} expert Delphi-style item rating & Canadian adult women, pooled multi-study sample & Developed a brief, validated 15-item short form of the Composite Abuse Scale, addressing the problem that most IPV \seqsplit{measurement} 'privileges physical violence, with poor assessment of other \seqsplit{experiences},' which had led to \seqsplit{underestimating} the scope of IPV in \seqsplit{populations}. (Verifier correction recorded in the library entry.) & Supports\\
\\[3pt]
Hegarty 2026 \cite{ref98} & instrument-\seqsplit{development}\slash{}validation study, n=854 adult women + \seqsplit{qualitative} survivor\slash{}expert feedback & Australian women, community survey sample & Developed and \seqsplit{psychometrically} tested a new self-report Coercive-Composite Abuse Scale (C-CAS) designed to capture severity, frequency and pattern of coercive control \seqsplit{holistically}, explicitly because prior measures showed 'limited \seqsplit{understanding} of these patterns of behaviour.' & Supports\\
\\[3pt]
Kurnaz Kaan 2026 \cite{srp42015126} & Cross-sectional \seqsplit{descriptive}-\seqsplit{correlational} survey with path analysis & 480 married women aged 18-49, Türkiye & Used the \seqsplit{Reproductive} Coercion Scale, Fertility Desire Scale, and Gender Roles Attitude Scale together; \seqsplit{reproductive} coercion and \seqsplit{traditional} gender-role attitudes were linked to fertility desire pathways in married women (path analysis). (Verifier correction recorded in the library entry.) & Supports\\
\\[3pt]
Kyegombe N 2018 \cite{srp30613427} & Cross-sectional survey nested within a cluster randomised trial (SASA! study) & Community members aged 18-49 who had witnessed IPV, Kampala, Uganda, 4 years post-\seqsplit{intervention} & SASA! community members who had witnessed IPV were \seqsplit{substantially} more likely to try to help than control-community \seqsplit{counterparts} (57\% vs 31\%); associated factors differed by trial arm, with SASA!-community bystander action linked to \seqsplit{relationship} duration, employment, talking to neighbours and SASA! exposure. & Supports\\
\\[3pt]
MacMillan 2006 \cite{ref96} & cluster randomized trial, n=2,602 eligible women across 6 clinical sites (emergency, family practice, women's health) & English-speaking women aged 18-64 in Ontario, Canada healthcare settings & Compared face-to-face interview, written self-completed \seqsplit{questionnaire}, and \seqsplit{computerized} screening for IPV to determine the optimal screening method in terms of accuracy, \seqsplit{acceptability} and \seqsplit{completeness}. (Verifier correction recorded in the library entry.) & Supports\\
\\[3pt]
O'Doherty L 2015 \cite{ref38} & Cochrane systematic review and meta-analysis of 13 RCTs\slash{}quasi-RCTs (n=14,959) & Women in antenatal, women's health, emergency department and primary care settings, mostly high-income urban settings & Screening increased clinical \seqsplit{identification} of IPV victims (OR 2.95, 95\% CI 1.79-4.87, n=10,074) but there was no evidence of effect on referrals to support services (OR 2.24, 95\% CI 0.64-7.86, only 2 studies, n=1298), no evidence screening reduced women's re-exposure to violence at 3-18 months, and \seqsplit{insufficient} evidence on health outcomes or harm beyond the immediate post-screening period; authors conclude screening alone is \seqsplit{insufficiently} supported to justify universal \seqsplit{implementation} absent stronger referral\slash{}outcome evidence. & Supports\\
\\[3pt]
Onukwugha FI 2025 \cite{ref42} & Quasi-\seqsplit{experimental} matched-pair cluster-controlled trial (n=1617 clients, 20 \seqsplit{intervention} vs 20 comparison clinics) & Family planning and antenatal care clients, Nigeria & Combined WHO LIVES + ARCHES first-line response reduced relative risk of IPV exposure (beta=0.54, 95\% CI 0.30-0.98) and \seqsplit{reproductive} coercion (beta=0.51, 95\% CI 0.28-0.94) over 9 months among women not already \seqsplit{experiencing} violence, but showed no effect on modern \seqsplit{contraceptive} use; authors note LIVES and ARCHES \seqsplit{individually} have shown 'mixed effects' in prior trials. (Verifier correction recorded in the library entry.) & Supports\\
\\[3pt]
Williamson HC 2014 \cite{ref106} & Cross-sectional survey study (N=2,126 married \seqsplit{individuals}) & US married adults, general population sample & Contrary to the assumption that premarital education reduces later need for counseling, receiving premarital education covaried with an INCREASED likelihood of later couples counseling (a 'gateway' effect), with the \seqsplit{association} stronger among African Americans and those with lower income\slash{}education. & Supports\\
\\[3pt]
Williamson HC 2019 \cite{ref107} & Cross-sectional self-report study (N=231 ethnically diverse newlywed couples in low-income \seqsplit{neighborhoods}) & US low-income newlywed couples & Top perceived barriers to seeking couple therapy were cost and \seqsplit{uncertainty} about where to go for help; direct experience with premarital education (or indirect experience via a social-network member's couple therapy) was associated with reduced barriers and greater likelihood of \seqsplit{subsequently} seeking therapy. & Supports\\
\\[3pt]
\seqsplit{Nsengiyumva} R 2025 \cite{srp41454335} & \seqsplit{Qualitative} \seqsplit{descriptive} study (focus group \seqsplit{discussions}, N=40 engaged couples) & Engaged couples recruited at premarital counseling sites (churches, sector \seqsplit{administration} offices), Rwanda & Engaged couples identified specific unmet \seqsplit{information}\slash{}education needs regarding conception \seqsplit{preparedness}, with premarital counseling sites (including religious venues) identified as a natural point of contact; themes included knowledge gaps, motivators for seeking \seqsplit{information}, and desired \seqsplit{educational} content. (Verifier correction recorded in the library entry.) & Supports\\
\\[3pt]
Evans 2016 \cite{ref51} & \seqsplit{qualitative} think-aloud\slash{}cognitive \seqsplit{interviewing} study with a validated \seqsplit{quantitative} scale (Composite Abuse Scale) & women with experience of domestic violence and abuse & Think-aloud \seqsplit{interviewing} revealed systematic \seqsplit{underreporting} on a widely used, validated DVA scale, \seqsplit{particularly} for coercive control, \seqsplit{threatening} behavior, \seqsplit{restrictions} to freedom, and sexual abuse, driven by item-\seqsplit{interpretation} confusion, fear of answering truthfully, and reluctance to self-identify as abused. & Complicates\\
\\[3pt]
Bhandari TR 2014 \cite{srp26982668} & Scale \seqsplit{construction} and validation (24-item pool, \seqsplit{psychometric} testing, cross-validation across two districts) & Women in Rupandehi and Kapilvastu districts, Nepal (South Asian, \seqsplit{patriarchal} marital-household context) & Produced a 23-item women's autonomy scale with strong \seqsplit{psychometric} properties (Cronbach's alpha 0.84, test-retest r=0.87, content validity ratio 0.8), validated \seqsplit{specifically} for South Asian household\slash{}marital decision contexts, showing good convergent and \seqsplit{discriminant} validity. (Verifier correction recorded in the library entry.) & Relevant only\\
\\[3pt]
World Health \seqsplit{Organization} 2013 (no PMID) & Global clinical and policy guideline & Global health-system guidance for responding to women subjected to IPV or sexual violence & \seqsplit{Establishes} WHO's first-line support principles for clinicians (the basis of the later-named LIVES model: Listen, Inquire about needs, Validate, Enhance safety, Support\slash{}refer), recommends against mandatory universal screening in the absence of adequate response systems, and sets minimum \seqsplit{requirements} (privacy, \seqsplit{confidentiality}, trained staff, referral pathways) before any clinical inquiry about violence should occur. & Relevant only\\
\\[3pt]
World Health \seqsplit{Organization} 2014 (no PMID) & Global clinical practice handbook \seqsplit{operationalizing} the 2013 WHO guidelines & Frontline health workers globally, first-line clinical response to women disclosing IPV or sexual violence & \seqsplit{Operationalizes} the LIVES first-line support protocol (Listen, Inquire, Validate, Enhance safety, Support) into a step-by-step clinical tool for frontline providers, with explicit guidance on \seqsplit{confidentiality}, \seqsplit{documentation}, safety planning, and structured referral rather than passive \seqsplit{information} provision. & Relevant only\\
\\[3pt]
\end{longtable}
